\documentclass[a4paper,10pt]{article}
\usepackage[utf8]{vietnam}
\usepackage[english]{babel}
\usepackage{amsmath}
\usepackage{amsfonts}
\usepackage{amssymb}
\usepackage{graphicx} % Required for inserting images
\usepackage{fancyhdr}
\usepackage{lastpage} 
\usepackage[dvipsnames]{xcolor}
\usepackage{tcolorbox}
\usepackage{titlesec}
\usepackage{mdframed}
\setlength{\parindent}{0pt} 
\setlength{\parskip}{0cm}
\usepackage{url}
\usepackage{pst-all}
\usepackage{pgfplots}
% \pgfplotsset{compat=1.15}
\usepackage{mathrsfs}
\usetikzlibrary{arrows} 
\usepackage{minted}
\usemintedstyle{emacs}
\usepackage{listings}
\usepackage{hanging}
\usepackage{subcaption}
\usepackage{float}
\usepackage{booktabs}
\usepackage{tabularx}
\usepackage[numbers,sort&compress]{natbib}
\pagestyle{empty}
\lstset{
    language=Python,              % Ngôn ngữ lập trình (Python)
    basicstyle=\ttfamily\small,   % Kiểu chữ cơ bản (monospace, cỡ nhỏ)
    keywordstyle=\color{green},    % Màu từ khóa
    stringstyle=\color{blue},      % Màu chuỗi
    commentstyle=\color{orange!90},    % Màu chú thích
    numbers=left,                 % Hiển thị số dòng (left hoặc right)
    numberstyle=\tiny\color{gray},% Kiểu chữ của số dòng
    stepnumber=1,                 % Hiển thị số dòng cách nhau 1 dòng
    frame=single,                 % Khung bao quanh mã (single/double)
    tabsize=4,                    % Kích thước của tab
    breaklines=true,              % Tự động xuống dòng
    captionpos=b                  % Đặt chú thích (t trên hoặc b dưới)
    literate={"}{{\texttt{"}}}1,
}
\usepackage{color}
% \usepackage{draftwatermark} 
\usepackage[thref,thmmarks,standard,amsmath]{ntheorem} 
\usepackage[colorlinks]{hyperref}
\hypersetup{
    colorlinks=true, 
    % linkcolor=blue,
    citecolor=blue,
    urlcolor=blue,
    filecolor=blue
}
%Môi trường Bài toán 
\theorembodyfont{\normalfont} 
\theoremseparator {.}   
\newtheorem{problem}{Bài toán}
%Môi trường Định lý
\theoremstyle{nonumberplain} 
\renewtheorem{theorem}{Định lý}
\theoremstyle{nonumberplain} 
\renewtheorem{corollary}{Hệ quả} 
%Môi trường bình luận 
\theoremstyle{nonumberplain} 
\newtheorem{cmt}{Bình luận} 
%Môi trường chú ý
\renewtheorem{remark}{Chú ý}
%Môi trường Bổ đề
%Môi trường Chứng minh
\theoremstyle{nonumberplain}  
\theoremseparator{.}
\theoremsymbol{\ensuremath{_\square}}
\renewtheorem{proof}{Lời giải}
\usepackage[left=1.5cm, right=1.5cm, top=1.5cm, bottom=1cm]{geometry} 
\newcommand{\thanh}[1]{{\textbf{\color{blue}{Thanh: #1}}}}

% \pagestyle{fancy}
% \fancyhf{} 
% \setlength{\footskip}{30pt}
% \fancyhead[L]{ \includegraphics[width=1.5cm]{VinUni_logo.png}
% }   % Left side
% \fancyhead[R]{College of Engineering and Computer Science \\ MATH4010 - Advanced Probability and Statistics}  % Right side

% \fancyfoot[L]{MATH4010-Spring 2026}   % Left side
% \fancyfoot[C]{} 
% \fancyfoot[R]{\thepage /\textcolor{blue}{6}}   
% \renewcommand{\footrulewidth}{0.4pt} 

% \fancypagestyle{firstpage}{
%     \fancyhf{} 
%     \fancyfoot[L]{}   % Left side
%     \fancyfoot[C]{} 
%     \fancyfoot[R]{} 
% }

\begin{document}
\textbf{MATH4010 Final Project: Group D}
\vspace{-0.2cm} 
\begin{table}[H]
\centering
\begin{tabular}{|c|l|c|l|}
\hline
\textbf{No.} & \textbf{Name} & \textbf{Student ID} & \textbf{Email} \\ \hline
1 & Nguyen Cung Thanh & V202401555 & 24thanh.nc@vinuni.edu.vn \\ \hline
2 & Duong Pham Minh Dung & V202401386 & 24dung.dpm@vinuni.edu.vn \\ \hline
3 & Nguyen Duy Minh & V202401435 &  24minh.nd@vinuni.edu.vn\\ \hline
4 & Do Chi Hai & V202401459  & 24hai.dc@vinuni.edu.vn  \\ \hline
5 & Nguyen Quang Huy & V202401400 &  24huy.nq@vinuni.edu.vn \\ \hline
6 & Nguyen Huy An Khanh & V202401646 & 
24khanh.nha@vinuni.edu.vn \\ \hline 
7 & Hoang Anh Tu & V202401388 & 24tu.ha@vinuni.edu.vn \\ \hline
8 & Kim Vladimir& V202401663&  24vladimir.k@vinuni.edu.vn\\ \hline 
9 & Nguyen Khanh Ngoc & V202401528 & 24ngoc.nk@vinuni.edu.vn  \\ \hline
10 & Ngo Thanh An & V202401748 & 24an.nt2@vinuni.edu.vn \\ \hline
\end{tabular}
\end{table}
\vspace{-0.3cm} 
Link to Github is \href{https://github.com/duongphamminhdung/MATH4010_Final_Project}{here}.
\vspace{-0.3cm} 


\section{Question 1}
\subsection{Question 1.a}

\textbf{Visualization.} In 2-dimensional space, we consider two circles: $A_2$ centered at (0.5,0.5) with radius 0.5, and $B_2$ centered at (0,0) with radius 1 (see Figure~\ref{a2} and Figure~\ref{b2}). 
\begin{figure}[hbt!]
    \centering
    \begin{minipage}{0.45\textwidth}
        \centering
        \begin{subfigure}{0.47\linewidth}
        \centering
        \includegraphics[scale = 0.27]{1.a/vis_A2.png}
        \caption{}
        \label{a2}
    \end{subfigure}
    \hfill
    \begin{subfigure}{0.5\linewidth}
        \centering
        \includegraphics[scale = 0.27]{1.a/vis_B2.png}
        \caption{}
        \label{b2}
    \end{subfigure}
    \end{minipage}%
    \hfill
    \begin{minipage}{0.5\textwidth}
        \begin{subfigure}{0.48\linewidth}
        \centering
        \includegraphics[scale = 0.20]{1.b/Maple(1).png}
        \caption{}
        \label{a3}
    \end{subfigure}
    \hfill
    \begin{subfigure}{0.47\linewidth}
        \centering
        \includegraphics[scale = 0.20]{1.b/Maple(2).png}
        \caption{}
        \label{b3}
    \end{subfigure}
    \end{minipage}
    \caption{Visualization of (a) $A_2$, (b) $B_2$ with their bounding boxes and (c) $A_3$, (d) $B_3$ with their bounding boxes.}
    \label{fig:visualization_1a}
\end{figure}
\vspace{-0.5cm}

\paragraph{Area calculation.} Using the standard formula for the area of a circle ($A = \pi r^2$), the areas of $A_2$ and $B_2$ are $\dfrac{\pi}{4}$ and $\pi$, respectively.  
\paragraph{Algorithm and Metrics.}
We estimate areas using the Monte Carlo hit-or-miss method. Given a shape $S$ inside a bounding region $\Omega$ with known volume $|\Omega|$, we draw $N$ uniform samples from $\Omega$ and count the hits $M = \#\{x_i \in S\}$. The area estimator is $\hat{A} = |\Omega| \cdot M/N$. To evaluate accuracy, we report three metrics: the relative error $|\hat{A} - A_{\text{true}}|/A_{\text{true}}$, the standard error $\text{SE} = |\Omega| \cdot \sqrt{\hat{p}(1-\hat{p})/N}$ where $\hat{p} = M/N$, and the 95\% confidence interval $\hat{A} \pm 1.96 \cdot \text{SE}$. The method generalizes directly to any dimension with convergence rate $O(1/\sqrt{N})$.

\paragraph{Estimates, Errors and Metrics.}
We ran the Monte Carlo hit-or-miss algorithm with $N = 10^7$ samples. The results are reported in Table \ref{tab:q1a}. Both estimates achieve relative errors below $0.01\%$, and the 95\% confidence intervals contain the true values, confirming that the estimator is unbiased.

\begin{table}[hbt!]
\centering
\caption{Monte Carlo estimation results for 2D ($N = 10^7$).}
\begin{tabular}{cccccc}
\hline
\textbf{Shape} & \textbf{Ground Truth} & \textbf{Estimate} & \textbf{Rel Error} & \textbf{Std Error} & \textbf{95\% CI} \\
\hline
$A_2$ & $\pi/4 = 0.785398$ & $0.785432$ & $0.0044\%$ & $0.000130$ & $[0.785178,\; 0.785687]$ \\
$B_2$ & $\pi = 3.141593$ & $3.141413$ & $0.0057\%$ & $0.000519$ & $[3.140395,\; 3.142431]$ \\
\hline
\end{tabular}
\label{tab:q1a} 
\end{table}
\subsection{Question 1.b}
\textbf{Visualization.} In 3-dimensional space, we consider two spheres: $A_3$ centered at (0.5,0.5,0.5) with radius 0.5, and $B_3$ centered at (0,0,0) with radius 1 (see Figure~\ref{a3} and Figure~\ref{b3}).

\paragraph{Volume calculation.} Using the volume formula for spheres in 3-dimensional space, the volumes of $A_3$ and $B_3$ are $V(A_3) = \dfrac{\pi}{6}$ and  $V(B_3) = \dfrac{4}{3}\pi.$

\paragraph{Estimates, Errors and Metrics.} 
We apply the same Monte Carlo hit-or-miss algorithm with $N = 10^7$ samples. The full results for 3D are reported in Table~\ref{tab:q1b}. All estimates achieve relative errors below $0.06\%$, and the 95\% confidence intervals contain the true values.

\begin{table}[H]
\centering
\caption{Monte Carlo estimation results for 3D ($N = 10^7$).}
\label{tab:q1b}
\begin{tabular}{cccccc}
\hline
\textbf{Shape} & \textbf{Ground Truth} & \textbf{Estimate} & \textbf{Rel Error} & \textbf{Std Error} & \textbf{95\% CI} \\
\hline
$A_3$ & $\pi/6 = 0.523599$ & $0.523333$ & $0.0508\%$ & $0.000158$ & $[0.523023,\; 0.523642]$ \\
$B_3$ & $4\pi/3 = 4.188790$ & $4.187672$ & $0.0267\%$ & $0.001264$ & $[4.185196,\; 4.190148]$ \\
\hline
\end{tabular}
\end{table}

\paragraph{Comparison with 2-dimensional space.}

Comparing the 2D and 3D rows in Table~\ref{tab:q1a} and Table~\ref{tab:q1b}, the relative error increases from 2D to 3D ($0.005\%$ to $0.04\%$ for $B$). The standard error scales as $\text{SE} = V_{\text{box}} \cdot \sqrt{p(1-p)/N}$ where $p$ is the hit probability. For $B_2$: $V_{\text{box}} = 4$, $p \approx 0.785$. For $B_3$: $V_{\text{box}} = 8$, $p \approx 0.524$. The larger bounding box and lower hit probability both contribute to higher standard error in 3D. The convergence rate remains $O(1/\sqrt{N})$ in both dimensions, but the constant factor grows with dimension.

\subsection{Question 1.c}

\paragraph{Exact Volume Formula and Estimation.}

The volume of an $n$-dimensional ball of radius $R$ is given by \cite[Exercise~4.29]{vershynin2018} as of an $n$-dimensional ball of radius $R$ as
$V_n(R) = \dfrac{\pi^{n/2}}{\Gamma\!\left(\frac{n}{2}+1\right)}R^n$ so $V(A_n) = \dfrac{\pi^{n/2}}{\Gamma\!\left(\frac{n}{2}+1\right)}\cdot \left( \dfrac{1}{2} \right)^n$ and $V(B_n) =  \dfrac{\pi^{n/2}}{\Gamma\!\left(\frac{n}{2}+1\right)}$,
where $\Gamma(x)$ is the Gamma function. For our shapes at $n=100$:
$$
V(A_{100}) = \frac{\pi^{50}}{\Gamma(51)} \cdot \left(\frac{1}{2}\right)^{100} \approx 1.868 \times 10^{-70}, \qquad
V(B_{100}) = \frac{\pi^{50}}{\Gamma(51)} \approx 2.368 \times 10^{-40}
$$

\paragraph{Naive Monte Carlo Failure.}

Applying the same hit-or-miss algorithm with $N = 10^7$ samples yields zero hits for both $A_{100}$ and $B_{100}$. The cause is geometric: in higher dimensions, the volume of a ball constitutes an exponentially vanishing fraction of its bounding box. For $A_{100}$, the hit probability is $p_A = V(A_{100}) \approx 1.87 \times 10^{-70}$; for $B_{100}$, after accounting for the box volume $2^{100}$, the probability is $p_B = V(B_{100})/2^{100} \approx 1.87 \times 10^{-70}$. The expected number of hits is $\mathbb{E}[M] = N \cdot p \approx 1.87 \times 10^{-63}$, so observing $M = 0$ is essentially certain. Achieving even one expected hit would require $N \approx 5 \times 10^{69}$ samples, which is computationally infeasible.

\paragraph{Variance Reduction Techniques.}

To mitigate the vanishing hit rate, we tested two classical variance reduction methods: stratified sampling via Latin Hypercube \cite{mckay1979} and antithetic variates \cite{hammersley1956}. At low dimensions ($n \le 8$), stratified sampling provides a $1.5$--$2\times$ reduction in relative error. Antithetic variates showed no improvement, as the ball indicator function is not monotone in the sampling coordinates. However, both methods still rely on uniform sampling and therefore cannot increase the hit rate. At $n = 100$, both produce zero hits, same as naive MC.

\paragraph{Key Insight: Thin-Shell Geometry.}

Since classical variance reduction cannot solve the problem, we investigate the root cause. By the concentration of the norm \cite{vershynin2018}, almost all volume of an $n$-dimensional ball concentrates in a thin shell near its surface -- over 99\% lies within distance $5/n$ from the boundary. At $n = 100$, this shell is just $0.05$ thick. Uniform sampling from the bounding box misses it entirely, explaining why all uniform-based methods fail.

\paragraph{Improvement by Importance Sampling.}

Since the core issue is that uniform proposals miss the thin shell, we replace them with a Gaussian proposal $q(x) = \mathcal{N}(c, \sigma^2 I_n)$ that naturally concentrates its mass on a shell of radius $\sigma\sqrt{n}$. Setting $\sigma = R/\sqrt{n}$ aligns this shell with the ball's surface, ensuring that a large fraction of proposed samples land inside the ball. The importance sampling estimator \cite{agapiou2017} is:
$$\hat{V} = \frac{1}{N} \sum_{i=1}^N \frac{\textbf{1}_{\{x_i \in S\}}}{q(x_i)}.$$

With $N = 10^7$ samples, importance sampling achieves sub-percent relative error for both $A_{100}$ and $B_{100}$, where all other methods produce zero hits. The results are summarized in Table~\ref{tab:q1c}.

\begin{table}[hbt!]
\centering
\caption{Volume estimation methods at $n=100$.}
\label{tab:q1c}
\renewcommand{\arraystretch}{1.5}
\small
\begin{tabular}{lcccc}
\hline
& \multicolumn{2}{c}{$A_{100}$} & \multicolumn{2}{c}{$B_{100}$} \\
\cmidrule(lr){2-3} \cmidrule(lr){4-5}
\textbf{Method} & \textbf{Estimate} & \textbf{Rel Err} & \textbf{Estimate} & \textbf{Rel Err} \\
\hline
Naive MC          & 0 (M=0)                  & 100\%  & 0 (M=0)                  & 100\% \\
Stratified (LHS)  & 0 (M=0)                  & 100\%  & 0 (M=0)                  & 100\% \\
Antithetic         & 0 (M=0)                  & 100\%  & 0 (M=0)                  & 100\% \\
Importance Samp.  & $1.867 \times 10^{-70}$  & 0.08\% & $2.367 \times 10^{-40}$  & 0.03\% \\
\hline
\end{tabular}
\end{table}

\paragraph{Conclusion.} 

The naive Monte Carlo algorithm works well for $n = 2$ and $n = 3$ but starts to fail at higher dimensions due to the exponentially vanishing hit rate. Variance reduction techniques offer only slight efficiency gains at low dimensions and do not increase the hit rate, so they cannot overcome the exponential geometry. Gaussian importance sampling with $\sigma = R/\sqrt{n}$ succeeds because it matches the thin-shell geometry of high-dimensional balls, achieving sub-0.1\% relative error at $n = 100$ with $N = 10^7$ samples. As $N \to \infty$, the importance sampling estimator converges at the standard Monte Carlo rate $O(1/\sqrt{N})$ by the central limit theorem, but with a variance constant governed by the $\chi^2$-divergence between the proposal and the target \cite{agapiou2017}. The Gaussian proposal $q(x) = \mathcal{N}(c, \sigma^2 I_n)$ with $\sigma = R/\sqrt{n}$ places its mass on the same thin shell where the ball's volume concentrates, keeping this divergence small even at $n = 100$. This explains why importance sampling succeeds where all uniform-based methods fail: it is not merely a variance reduction trick, but a change of measure that matches the geometry of high-dimensional balls.

\section{Question 2}

\paragraph{Normalization Factor.}

Given $S(x) = x^4$, the target distribution is $f(x) = \frac{1}{Z}e^{-x^4}$ where $Z$ is the normalization constant ensuring $f$ integrates to one. Since $e^{-x^4}$ is an even function, we write
\[
Z = 2\int_0^\infty e^{-x^4}\mathrm{d}x \overset{u=x^4}{=}
2\int_0^\infty e^{-u}\cdot \frac14 u^{1/4-1}\mathrm{d}u = 
\frac12\,\Gamma\!\left(\frac14\right)
\approx 1.8128.
\]

On the other hand, because $f(x) \propto e^{-x^4}$ is symmetric about $x = 0$, the mean is $\mathbb{E}[X] = 0$. Now, \[\mathbb E[X^2] = \frac1Z\int_{-\infty}^{\infty} x^2 e^{-x^4}\mathrm{d}x =
\frac{2}{Z}\int_0^\infty x^2 e^{-x^4}\mathrm{d}x \overset{u=x^4}{=} \frac{2}{Z}\cdot\frac14
\int_0^\infty u^{-1/4}e^{-u}\mathrm{d}u = \dfrac{1}{2Z}\Gamma\left({\frac34}\right) = \frac{\Gamma(3/4)}{\Gamma(1/4)}
\approx 0.3380, 
\]
which gives $\mathrm{Var}(X) = \mathbb{E}[X^2] - \mathbb{E}[X]^2 = \mathbb{E}[X^2] \approx 0.3380.$ 

% \paragraph{Metropolis Sampling.} 

% The Metropolis algorithm \cite{metropolis1953} constructs a Markov chain whose stationary distribution equals the target $\pi(x)$. At each step, the algorithm proposes a candidate $y$ from a symmetric proposal distribution $q(y|x) = q(x|y)$, computes the acceptance ratio $\alpha = \min(1, \pi(y)/\pi(x))$, and accepts $y$ with probability $\alpha$ or retains the current state. Because the proposal is symmetric, the Hastings correction $q(x|y)/q(y|x)$ equals one. We use the unnormalized density $\pi(x) = e^{-x^4}$, since $Z$ cancels in the ratio.

\paragraph{Target Scale and Proposal Tuning.}

The efficiency of the random-walk Metropolis algorithm \cite{metropolis1953} depends strongly on the proposal step size. In this problem, we use the unnormalized target density $\pi(x) = e^{-x^4}$, since the normalization constant $Z$ cancels in the Metropolis acceptance ratio. A proposal variance that is too small leads to slow exploration of the target distribution, while a variance that is too large causes frequent rejection of proposed moves. For this reason, the proposal scale is commonly chosen relative to the natural spread of the target distribution, measured by its standard deviation $\sigma_{\text{target}} = \sqrt{\mathrm{Var}(X)} \approx 0.5814.$

\paragraph{Proprosal Distribution and Optimal Scaling.} \cite{roberts1997} established the optimal scaling framework for random-walk Metropolis algorithms, proving that the proposal variance can be chosen to maximize convergence speed. \cite[Table~1]{gelman1996} report the numerically optimal acceptance rate for $d = 1$ as approximately 44\%, achieved at proposal scale $\sigma_{\text{opt}} = 2.38 \cdot \sigma_{\text{target}}$. For our target, this gives $\sigma_{\text{opt}} = 2.38 \times 0.5814 \approx 1.38$. We apply this step size uniformly to three symmetric proposal distributions: Gaussian $q(y|x) = \mathcal{N}(x, \sigma^2)$, Uniform $q(y|x) = \text{Uniform}(x - \sigma, x + \sigma)$, and Laplace $q(y|x) = \text{Laplace}(x, \sigma)$. Each chain runs for $N = 1000$ iterations starting at $x_0 = 0$ with a fixed seed for reproducibility.

\paragraph{Numerical Performance.} The results are shown in Table~\ref{tab:q2_initial} and Figure~\ref{fig:q2_initial}. The Gaussian proposal achieves an acceptance rate of 48.15\%, reasonably close to the theoretical target of 44\%. The Laplace proposal gives a comparable acceptance rate of 43.94\%, while the Uniform proposal produces a substantially higher acceptance rate of 63.06\%. All three proposals generate sample means close to zero and sample variances close to the theoretical value of 0.3380, indicating that the Markov chains successfully sample from the target distribution.
\begin{table}[hbt!]
\centering
\caption{Results with $\sigma_{\text{opt}} \approx 1.38$ applied to all proposals ($N = 1000$).}
\label{tab:q2_initial}
\begin{tabular}{lcccc}
\hline
\textbf{Proposal} & \textbf{Step Size} & \textbf{Acceptance Rate} & \textbf{Sample Mean} & \textbf{Sample Variance} \\
\hline
Gaussian & 1.3837 & 48.15\% & $-0.0213$ & 0.3745 \\
Uniform & 1.3837 & 63.06\% & $0.0276$ & 0.3249 \\
Laplace & 1.3837 & 43.94\% & $0.0168$ & 0.3455 \\
\hline
Theoretical & -- & $\sim$44\% & $0.0000$ & 0.3380 \\
\hline
\end{tabular}
\end{table}

\begin{figure}[hbt!]
    \centering
    \begin{subfigure}{0.32\linewidth}
        \centering
        \includegraphics[width=\linewidth]{2/fig_initial_gaussian.png}
        \caption{Gaussian}
    \end{subfigure}
    \hfill
    \begin{subfigure}{0.32\linewidth}
        \centering
        \includegraphics[width=\linewidth]{2/fig_initial_uniform.png}
        \caption{Uniform}
    \end{subfigure}
    \hfill
    \begin{subfigure}{0.32\linewidth}
        \centering
        \includegraphics[width=\linewidth]{2/fig_initial_laplace.png}
        \caption{Laplace}
    \end{subfigure}
    \caption{Density histograms for the three proposals with $\sigma_{\text{opt}} = 1.38$ applied to all. The red curve is the true density $f(x)/Z$.}
    \label{fig:q2_initial}
\end{figure}
The differences in acceptance rates arise from the interaction between the proposal distributions and the target density. The optimal scaling rule $\sigma_{\text{opt}} = 2.38\cdot\sigma_{\text{target}}$ was primarily developed for Gaussian random-walk proposals \cite{roberts1997}. Since the target density $f(x)\propto e^{-x^4}$ decays faster than a Gaussian density, different proposal shapes produce different exploration behavior at the same step size. The Uniform proposal generates bounded moves, leading to smaller effective jumps and therefore higher acceptance rates. In contrast, the Laplace proposal combines frequent local moves with occasional larger jumps due to its heavier tails. Overall, the Gaussian and Laplace proposals produce acceptance rates closest to the theoretical optimum while maintaining accurate estimation of the target moments.

\paragraph{Improvement by Step Size Tuning and Comparison.}

Since the theoretical scaling rule $\sigma_{\mathrm{opt}} = 2.38\cdot\sigma_{\text{target}}$ may not remain optimal for non-Gaussian proposal distributions, we additionally tune the step size for each proposal independently by minimizing $|\text{acceptance rate} - 0.44|$. The tuned step sizes and corresponding results are shown in Table~\ref{tab:q2_tuned} and Figure~\ref{fig:q2_tuned}.

\begin{table}[H]
\centering
\caption{Tuned results with per-proposal optimized step sizes ($N = 1000$).}
\label{tab:q2_tuned}
\begin{tabular}{lcccc}
\hline
\textbf{Proposal} & \textbf{Step Size} & \textbf{Acceptance Rate} & \textbf{Sample Mean} & \textbf{Sample Variance} \\
\hline
Gaussian & 1.6466 & 44.04\% & $-0.0086$ & 0.3748 \\
Uniform & 1.9716 & 46.05\% & $-0.0049$ & 0.3298 \\
Laplace & 1.3581 & 44.04\% & $0.0218$ & 0.3440 \\
\hline
Theoretical & -- & $\sim$44\% & $0.0000$ & 0.3380 \\
\hline
\end{tabular}
\end{table}

\begin{figure}[hbt!]
    \centering
    \begin{subfigure}{0.32\linewidth}
        \centering
        \includegraphics[width=\linewidth]{2/fig_tuned_gaussian.png}
        \caption{Gaussian}
    \end{subfigure}
    \hfill
    \begin{subfigure}{0.32\linewidth}
        \centering
        \includegraphics[width=\linewidth]{2/fig_tuned_uniform.png}
        \caption{Uniform}
    \end{subfigure}
    \hfill
    \begin{subfigure}{0.32\linewidth}
        \centering
        \includegraphics[width=\linewidth]{2/fig_tuned_laplace.png}
        \caption{Laplace}
    \end{subfigure}
    \caption{Density histograms after per-proposal step size tuning. All proposals achieve acceptance rates near 44\%.}
    \label{fig:q2_tuned}
\end{figure}

After tuning, the Gaussian and Laplace proposals both achieve acceptance rates of approximately $44\%$, while the Uniform proposal remains slightly higher at $46.05\%$. All three proposals continue to produce sample means close to zero and sample variances near the theoretical value of $0.3380$, indicating accurate sampling from the target distribution. The tuned step sizes differ from the theoretical $\sigma_{\text{opt}} = 1.38$: the Gaussian requires a larger step (1.65), the Uniform requires an even larger step (1.97), while the Laplace stays close to the theoretical value (1.36). This confirms that the $2.38 \cdot \sigma_{\text{target}}$ formula provides a reasonable starting point, but per-proposal optimization is necessary for best performance, particularly when the proposal distribution differs from the Gaussian assumed in the original derivation.

\section{Question 3} 
\subsection*{3.a. Data Collection and Preprocessing}
Ambient noise clips were recorded on an iPhone 15 Pro and music clips on an iPhone 13 Pro, both using the built-in Voice Memos app. The recordings were converted from their native format to WAV with the following unified specifications: mono channel, 16\,kHz sampling rate, 16-bit depth. Each clip is 4 seconds long.

Each clip was then preprocessed in two steps: (1) temporal trimming to remove boundary artifacts, and (2) DC offset removal by subtracting the sample mean.
\subsection*{3.b. Compute four numbers per clip.}
We select the bin width $\Delta$ using the Freedman-Diaconis rule \cite{an2025}, $\Delta = 2 \cdot \text{IQR}(x) / n^{1/3}$, where $x$ is the concatenated amplitude samples from all 20 clips and $n$ is the total number of samples. Applied to our 20 clips, this yields $\Delta \approx 0.000281$ with $7{,}122$ bins ($H_{\max} = 12.798$). The per-clip results, including the entropy gap $H_{\text{Gauss}} - \widehat{H}$, are summarized in Table~\ref{tab:q3_per_clip}.


% For each clip we then compute four quantities: the discrete entropy $\widehat{H}_Q$, the differential entropy estimate $\widehat{H} = \widehat{H}_Q + \log_2 \Delta$, the maximum entropy $H_{\max} = \log_2 K$ where $K$ is the number of non-empty bins, and the Gaussian reference entropy $H_{\text{Gauss}} = \frac{1}{2}\log_2(2\pi e \hat{\sigma}^2)$. 

\begin{table}[H]
    \centering
    \caption{Per-clip entropy metrics for ambient noise and music clips.}
    \label{tab:q3_per_clip}
    {\normalsize
    \setlength{\tabcolsep}{4pt}
    \begin{minipage}[t]{0.48\textwidth}
        \centering
        \textbf{Ambient}\\
        \vspace{0.1cm}
        \begin{tabular}{c c c c c c}
    \hline
    \textbf{No.} & $\mathbf{H_Q}$ & $\mathbf{\hat{H}}$ & $\mathbf{H_{max}}$ & $\mathbf{H_{gauss}}$ & \textbf{Gap} \\
    \hline
    1 & 9.176 & -2.622 & 12.798 & -2.555 & 0.067 \\
    2 & 8.705 & -3.093 & 12.798 & -3.007 & 0.086 \\
    3 & 7.868 & -3.930 & 12.798 & -3.923 & 0.007 \\
    4 & 5.436 & -6.361 & 12.798 & -6.359 & 0.002 \\
    5 & 5.329 & -6.469 & 12.798 & -6.467 & 0.001 \\
    6 & 5.192 & -6.606 & 12.798 & -6.602 & 0.004 \\
    7 & 9.020 & -2.778 & 12.798 & -2.749 & 0.029 \\
    8 & 5.985 & -5.812 & 12.798 & -5.809 & 0.003 \\
    9 & 4.919 & -6.879 & 12.798 & -6.874 & 0.005 \\
    10 & 8.760 & -3.038 & 12.798 & -3.015 & 0.023 \\
    \hline
    \textbf{mean} & 7.039 & -4.759 & 12.798 & -4.736 & \textbf{0.023} \\
    \textbf{std.} & 1.808 & 1.808 & 0.000 & 1.830 & 0.030 \\
    \hline
    \end{tabular}
    \end{minipage}\hfill
    \begin{minipage}[t]{0.48\textwidth}
        \centering
        \textbf{Music}\\
        \vspace{0.1cm}
        \begin{tabular}{c c c c c c}
            \hline
            \textbf{No.} & $\mathbf{H_Q}$ & $\mathbf{\hat{H}}$ & $\mathbf{H_{max}}$ & $\mathbf{H_{gauss}}$ & \textbf{Gap} \\
            \hline
            1 & 7.228 & -4.570 & 12.798 & -4.452 & 0.118 \\
            2 & 7.721 & -4.077 & 12.798 & -3.914 & 0.163 \\
            3 & 7.392 & -4.405 & 12.798 & -4.371 & 0.034 \\
            4 & 8.190 & -3.608 & 12.798 & -3.575 & 0.033 \\
            5 & 8.500 & -3.297 & 12.798 & -3.251 & 0.047 \\
            6 & 8.678 & -3.120 & 12.798 & -3.039 & 0.082 \\
            7 & 8.257 & -3.541 & 12.798 & -3.287 & 0.254 \\
            8 & 7.954 & -3.843 & 12.798 & -3.662 & 0.181 \\
            9 & 8.198 & -3.600 & 12.798 & -3.480 & 0.120 \\
            10 & 8.802 & -2.996 & 12.798 & -2.981 & 0.015 \\
            \hline
            \textbf{mean} & 8.092 & -3.706 & 12.798 & -3.601 & \textbf{0.105} \\
            \textbf{std.} & 0.522 & 0.522 & 0.000 & 0.512 & 0.077 \\
            \hline
        \end{tabular}
    \end{minipage}
    }
\end{table}

\subsection*{3.c. Discussion}

The entropy gap $H_{\text{Gauss}} - \widehat{H}$ measures how far the empirical
distribution deviates from a Gaussian with the same variance. By the
maximum-entropy property of the Gaussian \cite{cover2006}, for any continuous random variable $X$
with variance $\sigma^2$,
\[
  h(X) \;\leq\; \tfrac{1}{2}\log_2(2\pi e\,\sigma^2) \;=\; H_{\text{Gauss}},
\]
with equality if and only if $X$ is Gaussian. Consequently, the gap is non-negative, and smaller values indicate distributions closer to Gaussian.

\paragraph{Observation.}
The ambient noise clips have a mean entropy gap of $0.023$, compared with $0.105$ for the music clips (Table~\ref{tab:q3_per_clip}). Statistically, five ambient clips exhibit entropy gaps close to zero. However, this should not be interpreted as evidence that the underlying signals are truly Gaussian. These recordings correspond to the quietest clips in the dataset, whose amplitudes occupy only a small fraction of occupied bins after preprocessing and discretization. With such limited effective support, the histogram estimator cannot reliably resolve the underlying distribution shape, artificially reducing the estimated entropy gap.

\paragraph{Normalization.}
To resolve this issue, we peak-normalize all clips to $[-1, 1]$, which spreads quiet signals across the full bin range and allows the histogram estimator to recover the true distribution shape. After normalization, the Freedman-Diaconis rule yields $\Delta \approx 0.003692$ with $542$ bins ($H_{\max} = 9.082$). However, since $H_{\text{Gauss}} = \frac{1}{2}\log_2(2\pi e\,\hat{\sigma}^2)$ is entirely determined by variance, normalization removes natural loudness differences and can distort the gap-based ranking. We therefore also report the Shapiro-Wilk $W$ statistic \citep{shapiro1965}, a scale-invariant omnibus test for normality ($W = 1$ indicates perfect Gaussianity), which yields identical values before and after normalization. The normalized per-clip results are shown in Table~\ref{tab:q3_normalized}.

\begin{table}[H]
    \centering
    \caption{Per-clip entropy metrics after peak normalization with Shapiro-Wilk $W$.}
    \label{tab:q3_normalized}
    {\normalsize
    \begin{minipage}[t]{0.48\textwidth}
        \centering
        \textbf{Ambient}\\
        \vspace{0.1cm}
        \setlength{\tabcolsep}{3.5pt}
        \begin{tabular}{c c c c c c c}
    \hline
    \textbf{No.} & $\mathbf{H_Q}$ & $\mathbf{\hat{H}}$ & $\mathbf{H_{max}}$ & $\mathbf{H_{gauss}}$ & \textbf{Gap} & $\mathbf{W}$ \\
    \hline
    1 & 7.193 & -0.889 & 9.082 & -0.838 & 0.051 & 0.977 \\
    2 & 6.756 & -1.325 & 9.082 & -1.252 & 0.073 & 0.951 \\
    3 & 7.903 & -0.179 & 9.082 & -0.173 & 0.006 & 1.000 \\
    4 & 7.960 & -0.121 & 9.082 & -0.031 & 0.090 & 1.000 \\
    5 & 8.117 & 0.035 & 9.082 & 0.119 & 0.084 & 1.000 \\
    6 & 7.829 & -0.252 & 9.082 & -0.157 & 0.095 & 0.999 \\
    7 & 8.399 & 0.317 & 9.082 & 0.342 & 0.024 & 0.994 \\
    8 & 7.934 & -0.148 & 9.082 & -0.113 & 0.035 & 0.999 \\
    9 & 8.044 & -0.037 & 9.082 & -0.009 & 0.028 & 0.998 \\
    10 & 7.594 & -0.488 & 9.082 & -0.473 & 0.015 & 0.995 \\
    \hline
    \textbf{mean} & 7.773 & -0.309 & 9.082 & -0.258 & \textbf{0.050} & \textbf{0.991} \\
    \textbf{std.} & 0.479 & 0.479 & 0.000 & 0.474 & 0.033 & 0.016 \\
    \hline
    \end{tabular}
    \end{minipage}\hfill
    \begin{minipage}[t]{0.48\textwidth}
        \centering
        \textbf{Music}\\
        \vspace{0.1cm}
        \setlength{\tabcolsep}{3pt}
        \begin{tabular}{c c c c c c c}
            \hline
            \textbf{No.} & $\mathbf{H_Q}$ & $\mathbf{\hat{H}}$ & $\mathbf{H_{max}}$ & $\mathbf{H_{gauss}}$ & \textbf{Gap} & $\mathbf{W}$ \\
            \hline
            1 & 7.206 & -0.875 & 9.082 & -0.756 & 0.119 & 0.963 \\
            2 & 6.484 & -1.597 & 9.082 & -1.441 & 0.156 & 0.938 \\
            3 & 7.587 & -0.494 & 9.082 & -0.459 & 0.035 & 0.989 \\
            4 & 7.416 & -0.665 & 9.082 & -0.637 & 0.028 & 0.990 \\
            5 & 7.208 & -0.873 & 9.082 & -0.835 & 0.039 & 0.990 \\
            6 & 7.475 & -0.606 & 9.082 & -0.533 & 0.073 & 0.971 \\
            7 & 7.005 & -1.076 & 9.082 & -0.830 & 0.246 & 0.933 \\
            8 & 7.010 & -1.072 & 9.082 & -0.897 & 0.175 & 0.935 \\
            9 & 7.387 & -0.694 & 9.082 & -0.580 & 0.114 & 0.954 \\
            10 & 7.796 & -0.285 & 9.082 & -0.277 & 0.008 & 0.999 \\
            \hline
            \textbf{mean} & 7.257 & -0.824 & 9.082 & -0.724 & \textbf{0.099} & \textbf{0.966} \\
            \textbf{std.} & 0.368 & 0.368 & 0.000 & 0.317 & 0.077 & 0.025 \\
            \hline
        \end{tabular}
    \end{minipage}
    }
\end{table}

\paragraph{Conclusion.}
After peak normalization, the ambient mean gap ($0.050$) remains smaller than the music mean gap ($0.099$), confirmed by the Shapiro-Wilk statistic (ambient $W = 0.991$ vs.\ music $W = 0.966$). The gap, also known as negentropy \citep{cover2006}, measures how much entropy is lost relative to the maximum-entropy Gaussian with the same variance. Ambient noise's near-zero gap means its amplitude distribution is nearly as random as the variance permits. This is consistent with the central limit theorem, since ambient sound is a superposition of many weakly dependent sources converging to the maximum-entropy Gaussian. Music, with its deliberate harmonic and temporal structure, carries additional constraints that reduce entropy below this maximum.

\section{References}
\renewcommand{\refname}{}
\vspace{-1cm}

\begin{thebibliography}{9}

\bibitem[An et~al.(2025)]{an2025}
An, J. E., Tung, C., Do, C., \& Chen, W. (2025).
\textit{Optimal binning for Small-Angle neutron scattering data using the Freedman-Diaconis rule}.
Cornell University.

\bibitem[Agapiou et~al.(2017)]{agapiou2017}
Agapiou, S., Papaspiliopoulos, O., Sanz-Alonso, D., \& Stuart, A. M. (2017).
Importance sampling: Intrinsic dimension and computational cost.
\textit{Statistical Science}, 32(3), 405--431.

\bibitem[Cover \& Thomas(2006)]{cover2006}
Cover, T. M., \& Thomas, J. A. (2006).
\textit{Elements of information theory} (2nd ed.). Wiley.

\bibitem[Gelman et~al.(1996)]{gelman1996}
Gelman, A., Roberts, G. O., \& Gilks, W. R. (1996).
Efficient Metropolis jumping rules.
In J. M. Bernardo et al.\ (Eds.), \textit{Bayesian Statistics 5} (pp.~599--607). Oxford University Press.

\bibitem[Hammersley \& Morton(1956)]{hammersley1956}
Hammersley, J. M., \& Morton, K. W. (1956).
A new Monte Carlo technique: Antithetic variates.
\textit{Mathematical Proceedings of the Cambridge Philosophical Society}, 52(3), 449--475.

\bibitem[McKay et~al.(1979)]{mckay1979}
McKay, M. D., Beckman, R. J., \& Conover, W. J. (1979).
A comparison of three methods for selecting values of input variables in the analysis of output from a computer code.
\textit{Technometrics}, 21(2), 239--245.

\bibitem[Metropolis et~al.(1953)]{metropolis1953}
Metropolis, N., Rosenbluth, A. W., Rosenbluth, M. N., Teller, A. H., \& Teller, E. (1953).
Equation of state calculations by fast computing machines.
\textit{The Journal of Chemical Physics}, 21(6), 1087--1092.


\bibitem[Roberts et~al.(1997)]{roberts1997}
Roberts, G. O., Gelman, A., \& Gilks, W. R. (1997).
Weak convergence and optimal scaling of random walk Metropolis algorithms.
\textit{The Annals of Applied Probability}, 7(1), 110--120.

\bibitem[Shapiro \& Wilk(1965)]{shapiro1965}
Shapiro, S.~S., \& Wilk, M.~B. (1965).
An analysis of variance test for normality (complete samples).
\textit{Biometrika}, 52(3/4), 591--611.

\bibitem[Vershynin(2018)]{vershynin2018}
Vershynin, R. (2018).
\textit{High-dimensional probability: An introduction with applications in data science}.
Cambridge University Press.

\end{thebibliography}

\section{Contribution declaration}

\renewcommand{\arraystretch}{1.15}
\vspace{-0.5cm}

\begin{table}[h]
\centering
\resizebox{\textwidth}{!}{
\begin{tabular}{|p{8cm}|c|c|c|c|c|c|c|c|c|c|}
\hline
\textbf{Contribution} &
\textbf{Thanh} &
\textbf{Dung} &
\textbf{Minh} &
\textbf{Hai} &
\textbf{Huy} &
\textbf{Khanh} &
\textbf{Tu} &
\textbf{Kim} &
\textbf{Ngoc} &
\textbf{An} \\ \hline

\multicolumn{11}{|c|}{\textbf{Question 1}} \\ \hline

1a. Visualization and compute exact area of $A_2,B_2$ &
 & & & \checkmark\checkmark &
 & & \checkmark &
 & & \checkmark \\ \hline

1a. Implement Monte Carlo to estimate areas $A_2,B_2$ &
 & & & & & & \checkmark\checkmark &
 & & \checkmark \\ \hline

1b. Visualization and compute exact area of  $A_3,B_3$ &
\checkmark & & \checkmark\checkmark & &
\checkmark & & & & &
\\ \hline

1b. Implement Monte Carlo to estimate areas $A_3,B_3$ &
& & \checkmark & & \checkmark\checkmark & & &
 & & \\ \hline

1b. Report estimates, errors, and other metrics in 1a., 1b. to make a comparison &
& & & & & & & \checkmark\checkmark
 & & \\ \hline

1c. Volume formula derivation and discuss on Variance Reduction and Importance Sampling method &
\checkmark\checkmark &
\checkmark & & & & & & & \checkmark & \\ \hline

\multicolumn{11}{|c|}{\textbf{Question 2}} \\ \hline

Compute normalization constant &
\checkmark\checkmark & &
& & & & & & & \\ \hline

Proposal distributions
& & & \checkmark & & & & \checkmark & & &
\checkmark\checkmark \\ \hline

Performance analysis and proposed tuning method & & \checkmark\checkmark & & &
& & & & & \\ \hline

\multicolumn{11}{|c|}{\textbf{Question 3}} \\ \hline

3a. Collect data and data preprocessing & & & &
\checkmark & & & & \checkmark &
\checkmark\checkmark &
\\ \hline

3b. Finding bin width and entropy computation & & & & &
\checkmark & & &
\checkmark\checkmark &
\checkmark &
 \\ \hline

3c. Gaussianity discussion & & \checkmark & & & & \checkmark &
& & \checkmark\checkmark &
\\ \hline

\multicolumn{11}{|c|}{\textbf{Presentation and Report}} \\ \hline

Quality check and writing final report &
\checkmark\checkmark &
\checkmark & & & & & & & & \checkmark \\ \hline

Prepare presentation slides  & & & \checkmark\checkmark & &
\checkmark & &
\checkmark & & &
\\ \hline

Prepare presentation poster  & & & & \checkmark & &
\checkmark\checkmark & & & \checkmark &
\\ \hline

\end{tabular}
}
\caption{Contribution summary. $\checkmark$ indicate contributors, $\checkmark\checkmark$ indicates student in charge.}
\end{table}

\begin{table}[h]
\centering
\renewcommand{\arraystretch}{1.2}
\resizebox{\textwidth}{!}{
\begin{tabular}{|c|l|p{6cm}|c|c|c|c|c|c|c|c|c|c|}
\hline
\textbf{Task} & \textbf{Part} & \textbf{Sub-task} & \textbf{Thanh} & \textbf{Dung} & \textbf{Minh} & \textbf{Hai} & \textbf{Huy} & \textbf{Khanh} & \textbf{Tu} & \textbf{Kim} & \textbf{Ngoc} & \textbf{An} \\ \hline

1 & 1.a & Visualization of $A_2$, $B_2$ &  &  &  &  \checkmark \checkmark &  &  &  &  &  &  \\ \hline

1 & 1.a & Exact volume calculation of $A_2$, $B_2$ &  &  &  & &  &  &  &  &  & \checkmark\checkmark \\ \hline

1 & 1.a & Monte Carlo estimation for $A_2$, $B_2$ &  &  &  & \checkmark &  &  & \checkmark\checkmark &  &  & \checkmark \\ \hline

1 & 1.b & Visualize $A_3$, $B_3$ spheres & \checkmark &  & \checkmark \checkmark &  &  &  &  &  &  &  \\ \hline

1 & 1.b & Exact volume calculation of $A_3$, $B_3$ &  &  & \checkmark\checkmark &  &  &  &  &  &  & \\ \hline

1 & 1.b & Monte Carlo estimation for $A_3$, $B_3$ &  &  &  &  & \checkmark\checkmark &  &  &  &  & \\ \hline

1 & 1.b & Volume estimation, error and convergence analysis &  &  & \checkmark &  &  &  & & \checkmark\checkmark &  &  \\ \hline

1 & 1.b & Make comparison with 1.a &  &  & &  &  &  & & \checkmark\checkmark &  &  \\ \hline

1 & 1.c & Volume formula derivation and explanation of algorithm failure & \checkmark\checkmark &  &  &  &  &  &  &  & \checkmark &  \\ \hline

1 & 1.c & Analyze MC weakness and variance reduction &  & \checkmark\checkmark &  &  &  &  &  &  &  &  \\ \hline

1 & 1.c & Discuss Importance Sampling & \checkmark & \checkmark\checkmark &  &  &  &  &  &  &  &  \\ \hline

2 & 2.a & Estimate normalization constant $Z$ & \checkmark\checkmark &  &  &  &  &  &  &  &  &  \\ \hline

2 & 2.b & Uniform distribution implementation &  &  & \checkmark\checkmark &  &  &  &  &  &  &  \\ \hline

2 & 2.b & Laplace distribution implementation   &  &  &  &  &  &  &  &  &  & \checkmark\checkmark \\ \hline

2 & 2.b & Gaussian distribution implementation   &  &  &  &  &  &  & \checkmark\checkmark &  &  &  \\ \hline

2 & 2.c & Analyze metrics and weaknesses of proposals &  & \checkmark\checkmark &  &  &  &  &  &  &  &  \\ \hline

2 & 2.c & Optimize acceptance rate and step size tuning &  & \checkmark\checkmark &  &  &  &  &  &  &  &  \\ \hline

3 & 3.a & Gather and preprocess music data &  &  &  &  &  &  &  &  & \checkmark\checkmark &  \\ \hline

3 & 3.a & Gather and preprocess ambient data &  &  &  & \checkmark\checkmark &  &  &  &  &  &  \\ \hline

3 & 3.b & Compute entropy-related values per clip &  &  &  &  & \checkmark &  &  & \checkmark\checkmark &  &  \\ \hline

3 & 3.b & Visualize entropy gap distribution &  &  &  &  &  & \checkmark\checkmark &  &  &  &  \\ \hline

3 & 3.b & Improve entropy calculation method &  &  &  &  &  &  &  &  & \checkmark\checkmark &  \\ \hline

3 & 3.b & Decide bandwidth and calculate the discrete entropy
 &  &  &  &  & \checkmark\checkmark &  &  &  &  &  \\ \hline

3 & 3.c & Gaussianity discussion and explanations & &  &  &  & \checkmark  & \checkmark &  &  & \checkmark\checkmark &  \\ \hline

4 & Report & Finalize the report content & \checkmark\checkmark & \checkmark &  &  &  &  &  &  &  &  \\ \hline

4 & Slides & Prepare presentation slides &  &  & \checkmark &  & \checkmark  &  & \checkmark\checkmark &  &  & \\ \hline

4 & Poster & Prepare presentation poster &  &  &  & \checkmark \checkmark &  & \checkmark &  &  & \checkmark & \\ \hline
\end{tabular}
}

\caption{Contribution Declaration}
\label{tab:contribution_declaration}
\end{table}


\end{document} 
